\section{Энергия}
\subsection{Работа}
%%%%%%%%%%%%%%%%%%%%%%%%%%%%%%%%%%%%%%%%%%%%%%%%%%%%
%%%%%%%%%%%%%%%%%%%%%%%%%%%%%%%%%%%%%%%%%%%%%%%%%%%%
 \z{Какую минимальную работу надо совершить, чтобы из колодца глубины $H = 10$ м поднять на тросе ведро с водой массой $m = 8$ кг? Линейная плотность троса $\mu = 0.4$ кг/м.}
%
{$A = \left(mg + \frac{\mu g H}{2}\right)H=1000\,\text{Дж}$}
%%%%%%%%%%%%%%%%%%%%%%%%%%%%%%%%%%%%%%%%%%%%%%%%%%%%
\z[{\begin{wrapfigure}[5]{r}{0.3\textwidth}%
		\vspace{-12pt}%
		\begin{flushright}%
 			\includegraphics[width=0.3\textwidth]{7.1.2}%
		\end{flushright}%
	\end{wrapfigure}%
	}]%
	{Сила, действующая на снаряд массы $m$ в стволе орудия, нарастает равномерно от нуля до $F_0$ на участке ствола длины $l_1$, не меняется на участке ствола длины $l_2$ и, наконец, равномерно уменьшается до нуля на участке ствола длины $l_3$. Какова скорость снаряда при вылете?}
%
{$v = \sqrt{\frac{F_0}{m}(l_1 + 2l_2 + l_3)}$}
%%%%%%%%%%%%%%%%%%%%%%%%%%%%%%%%%%%%%%%%%%%%%%%%%%%% 
 \z{Груз, подвешенный на лёгкой пружине жёсткостью $k = 200$ Н/м, растягивает её на $x = 2$ см. Какую работу необходимо совершить вертикальной силе, приложенной к грузу, чтобы деформация пружины стала вдвое больше начальной?}
%
{$A=\frac{kx^2}{2}=0.04\, \text{Дж}$}
%%%%%%%%%%%%%%%%%%%%%%%%%%%%%%%%%%%%%%%%%%%%%%%%%%%% 
 \z{Транспортер поднимает массу $m = 200$ кг песка на автомашину за время $\tau = 1$ с. Длина ленты транспортера $L = 3$ м, угол её наклона к горизонту $\alpha = 30^\circ$. КПД транспортера $\eta = 0.85$. Найдите мощность $W$, развиваемую его электродвигателем.}
%
{$W = \frac{mgL\sin\alpha}{\eta\tau} \approx 3460\, \text{Вт}$}
%%%%%%%%%%%%%%%%%%%%%%%%%%%%%%%%%%%%%%%%%%%%%%%%%%%%
%%%%%%%%%%%%%%%%%%%%%%%%%%%%%%%%%%%%%%%%%%%%%%%%%%%%	
 \z[{\begin{wrapfigure}[4]{r}{0.3\textwidth}%
		\vspace{-12pt}%
		\begin{flushright}%
 			\includegraphics[width=0.3\textwidth]{7.1.5}%
		\end{flushright}%
	\end{wrapfigure}%
	}]%
	{Вверх по наклонной плоскости равномерно со скоростью $v$ поднимают тело массы $m$, причём сила направлена вдоль наклонной плоскости. При каком угле наклона $\alpha$ затрачиваемая мощность $W$ будет максимальной, и каково значение максимальной мощности? Коэффициент трения между телом и наклонной плоскостью $\mu = 1$.}
 %
{$\alpha_{\text{опт}} = 45\degree, ~W_{\text{max}} = mgv\sqrt{2}$}